\documentclass[11pt]{article}
\usepackage[margin=1in]{geometry}
\usepackage{enumitem}
% =============================================================================
% Preambles/header.tex — shared LaTeX/Beamer preamble for this project
%
% Use from a Beamer lecture by prepending:
%     \input{header}
% and compiling with TEXINPUTS=../Preambles:$TEXINPUTS (the /compile-latex
% skill does this automatically).
%
% PALETTE CONTRACT. Color names defined here MUST match the names in
% Quarto/theme-template.scss so Beamer and Quarto renderings use the same
% palette. The scripts/check-palette-sync.sh script enforces this.
%
% To customize: edit the HEX values below AND the matching $variable in
% Quarto/theme-template.scss, then run scripts/check-palette-sync.sh.
% =============================================================================

% -----------------------------------------------------------------------------
% Engine detection + encoding
%
% Our /compile-latex skill uses XeLaTeX, where inputenc is unnecessary and
% can produce encoding warnings. Only load inputenc under pdfTeX.
% -----------------------------------------------------------------------------
\usepackage{iftex}
\ifPDFTeX
  \usepackage[utf8]{inputenc}
\fi

\usepackage{amsmath, amssymb, amsthm}
\usepackage{booktabs}
\usepackage{graphicx}

% -----------------------------------------------------------------------------
% PALETTE — mirrors Quarto/theme-template.scss variable names.
% Load xcolor and define colors BEFORE anything that references them
% (notably hyperref, hypersetup, and the Beamer color assignments).
% -----------------------------------------------------------------------------
\usepackage{xcolor}

\definecolor{primary-blue}{HTML}{012169}      % headings, accents
\definecolor{primary-gold}{HTML}{B9975B}      % emphasis, borders
\definecolor{highlight-yellow}{HTML}{F2A900}  % markers, alerts
\definecolor{light-bg}{HTML}{E8EDF5}          % title / section backgrounds
\definecolor{jet}{HTML}{1A1A1A}               % body text

% Semantic colors (used in diagrams and annotations)
\definecolor{positive}{HTML}{15803D}          % good / observed / identified
\definecolor{negative}{HTML}{B91C1C}          % bad / problematic / bias
\definecolor{neutral}{HTML}{525252}           % reference / context

% Highlight accents (match .hi-* classes in the SCSS)
\definecolor{hi-slate}{HTML}{314F4F}
\definecolor{hi-green}{HTML}{2E8B57}
\definecolor{hi-red}{HTML}{C41E3A}

% Semantic aliases so skill templates and snippets can write
% \textcolor{accent}{...} without hard-coding "primary-blue".
\colorlet{accent}{primary-blue}
\colorlet{accent2}{primary-gold}

% -----------------------------------------------------------------------------
% Hyperref — loaded AFTER xcolor + palette so link colors resolve.
% -----------------------------------------------------------------------------
\usepackage{hyperref}
\hypersetup{colorlinks=true, linkcolor=primary-blue, urlcolor=primary-blue,
            citecolor=primary-blue, pdfborder={0 0 0}}

% -----------------------------------------------------------------------------
% Course code — set once per deck (after \input{header}, before \begin{document})
% via \coursecode{CS401}. Shown in the footer on every slide so a deck's course
% is identifiable without opening the title page. Empty by default.
% -----------------------------------------------------------------------------
\makeatletter
\newcommand{\coursecode}[1]{\gdef\@coursecodetext{#1}}
\gdef\@coursecodetext{}
\makeatother

% -----------------------------------------------------------------------------
% Beamer theme assignments (only apply under Beamer).
%
% We detect Beamer via \@ifclassloaded{beamer}, which is the documented,
% non-internal way. This must be wrapped in \makeatletter/\makeatother
% because \@ifclassloaded contains an @-catcode character.
% -----------------------------------------------------------------------------
\makeatletter
\@ifclassloaded{beamer}{%
  \usefonttheme{professionalfonts}
  \setbeamercolor{structure}{fg=primary-blue}
  \setbeamercolor{frametitle}{fg=primary-blue}
  \setbeamercolor{title}{fg=primary-blue}
  \setbeamercolor{subtitle}{fg=primary-gold}
  \setbeamercolor{author}{fg=primary-blue}
  \setbeamercolor{institute}{fg=neutral}
  \setbeamercolor{date}{fg=primary-gold}
  \setbeamercolor{itemize item}{fg=primary-blue}
  \setbeamercolor{itemize subitem}{fg=primary-gold}
  \setbeamercolor{alerted text}{fg=primary-gold}
  \setbeamercolor{block title}{fg=white, bg=primary-blue}
  \setbeamercolor{block body}{bg=light-bg}
  % Semantic block variants: block=definition/concept (blue), exampleblock=
  % worked example (green/positive), alertblock=key takeaway or caution (gold).
  % Keep to <=2 colored boxes per slide across ALL three variants (INV-7).
  \setbeamercolor{block title example}{fg=white, bg=positive}
  \setbeamercolor{block body example}{bg=light-bg}
  \setbeamercolor{block title alerted}{fg=white, bg=primary-gold}
  \setbeamercolor{block body alerted}{bg=light-bg}

  % Minimal footer: course code (left) + page X / N (right), no navigation clutter
  \setbeamertemplate{navigation symbols}{}
  \setbeamertemplate{footline}{%
    \color{neutral}\footnotesize\hspace{0.7em}\@coursecodetext\hfill
    \insertframenumber/\inserttotalframenumber\hspace{0.7em}\vspace{0.3em}}
}{}
\makeatother

% -----------------------------------------------------------------------------
% TikZ setup — libraries the snippet gallery uses
% -----------------------------------------------------------------------------
\usepackage{tikz}
\usetikzlibrary{arrows.meta, positioning, calc, decorations.pathreplacing,
                fit, shapes.geometric, backgrounds}

% Shared TikZ styles — snippets and hand-written diagrams can pick these up
% via `\begin{tikzpicture}[every node/.style={...}]` or by naming specific
% styles (e.g., `\node[dag-node] (X) at (0,0) {X};`).
\tikzset{
  % Node styles for causal DAGs and similar
  dag-node/.style = {
    draw=primary-blue, fill=light-bg, rounded corners=2pt,
    minimum width=1.2cm, minimum height=0.9cm, align=center, font=\small
  },
  decision-node/.style = {
    draw=primary-gold, fill=white, diamond, aspect=1.8,
    minimum width=1.8cm, minimum height=1.2cm, align=center, font=\small,
    inner sep=1pt
  },
  % Edge styles — observed vs counterfactual
  observed-edge/.style = {-{Latex[length=2.5mm]}, thick, draw=jet},
  counterfactual-edge/.style = {-{Latex[length=2.5mm]}, thick, draw=neutral, dashed},
  confound-edge/.style = {-{Latex[length=2.5mm]}, thick, draw=negative, dashed},
  % Data-point styles for plots
  observed-dot/.style = {fill=primary-blue, draw=primary-blue, circle, inner sep=1.2pt},
  counterfactual-dot/.style = {fill=white, draw=neutral, circle, inner sep=1.2pt}
}

% -----------------------------------------------------------------------------
% Convenience macros
% -----------------------------------------------------------------------------
\newcommand{\muted}[1]{\textcolor{neutral}{#1}}
\newcommand{\key}[1]{\textcolor{primary-gold}{\textbf{#1}}}
\newcommand{\good}[1]{\textcolor{positive}{#1}}
\newcommand{\bad}[1]{\textcolor{negative}{#1}}

% Standout frame macro: full-bleed dark background for section transitions.
% Beamer-only (uses \begin{frame}/\setbeamercolor/\usebeamerfont) -- do not
% call from an article-class file (e.g. Notes/<CODE>/*-notes.tex). \muted,
% \key, \good, \bad, and \coursecode above are plain xcolor/gdef and are
% safe to use from either class; verified when Notes/article-class support
% was added.
% Expand: \transitionslide{Your transition title}
\newcommand{\transitionslide}[1]{%
  \begingroup
  \setbeamercolor{background canvas}{bg=primary-blue}
  \begin{frame}[plain]
    \centering
    \vfill
    {\usebeamerfont{title}\color{white}\Large #1\par}
    \vfill
  \end{frame}
  \endgroup
}

% Section divider frame (Beamer-only), an alternative to \transitionslide with a
% darker two-line style: near-black (jet) background, a small white label line
% above a larger gold title, and a thin gold rule along the bottom edge.
% Expand: \sectiondivider{Section 2}{Pointers}
\newcommand{\sectiondivider}[2]{%
  \begingroup
  \setbeamercolor{background canvas}{bg=jet}
  \begin{frame}[plain]
    \vspace*{3.0cm}
    {\color{white}\ttfamily\large #1\par}
    \vspace{0.5cm}
    {\color{primary-gold}\LARGE #2\par}
    \begin{tikzpicture}[remember picture, overlay]
      \draw[primary-gold, line width=2.5pt]
        ($(current page.south west)+(0,0.1)$) -- ($(current page.south east)+(0,0.1)$);
    \end{tikzpicture}
  \end{frame}
  \endgroup
}

% -----------------------------------------------------------------------------
% Channel watermark — "Concepts That Click" (YouTube).
%
% Article-class files (CompetitiveExam Q/A sets, Notes, Assignments) get a
% light diagonal draft watermark on every page plus a centered footer naming
% the channel. Beamer decks get a subtle TikZ background watermark instead
% (the footline already carries the course code). Centralized here so every
% MCQ PDF is branded without per-file edits.
% -----------------------------------------------------------------------------
\makeatletter
\@ifclassloaded{article}{%
  \usepackage{draftwatermark}
  \SetWatermarkText{Concepts That Click}
  \SetWatermarkScale{1}
  \SetWatermarkFontSize{2cm}
  \SetWatermarkLightness{0.86}
  \SetWatermarkAngle{45}
  \usepackage{fancyhdr}
  \pagestyle{fancy}
  \fancyhf{}
  \fancyfoot[C]{\footnotesize\textcolor{neutral}{Concepts That Click~|~YouTube: \href{https://www.youtube.com/@concepts.thatclick}{@concepts.thatclick}}}
  \renewcommand{\headrulewidth}{0pt}
  \renewcommand{\footrulewidth}{0pt}
  \fancypagestyle{plain}{%
    \fancyhf{}%
    \fancyfoot[C]{\footnotesize\textcolor{neutral}{Concepts That Click~|~YouTube: \href{https://www.youtube.com/@concepts.thatclick}{@concepts.thatclick}}}%
    \renewcommand{\headrulewidth}{0pt}%
    \renewcommand{\footrulewidth}{0pt}%
  }
}{}
\@ifclassloaded{beamer}{%
  \setbeamertemplate{background}{%
    \begin{tikzpicture}[remember picture, overlay]
      \node[rotate=45, opacity=0.055, align=center]
        at (current page.center)
        {\Huge\textcolor{neutral}{Concepts That Click}};
    \end{tikzpicture}%
  }
}{}
\makeatother 
\coursecode{JKSSB}
\renewcommand{\thesection}{12.\arabic{section}}
\newtheorem{example}{Example}[section]
\newenvironment{definitionbox}[1]{%
  \par\noindent\textbf{Definition (#1).}\ \itshape
}{\par\vspace{0.3em}}
\title{Input Devices --- Detailed Lecture Notes}
\author{JKSSB Computer Awareness}
\date{\today}

\begin{document}
\maketitle

\section{What is an input device}
A computer works in a cycle: input, process, and output. In the input step,
data and instructions are fed into the machine. An \textbf{input device} is
the hardware that accepts data or control signals from the user or from the
environment and passes them to the computer, converting a human action or a
physical signal into machine-readable form. Without an input device the
computer would have nothing to process.

\begin{definitionbox}{Input device}
Hardware that accepts data or control signals from the user or the
environment and supplies them to the CPU, converting them into a form the
computer can process.
\end{definitionbox}

Common input devices include the keyboard, mouse, touchscreen, scanner,
microphone, digital camera, barcode reader, and biometric sensor. The
following table fixes the meaning of each term used in this lesson.

\begin{center}
\begin{tabular}{p{0.26\textwidth}p{0.66\textwidth}}
\toprule
\textbf{Term} & \textbf{Meaning in this lesson} \\
\midrule
Keyboard & Text input device that converts a keystroke into a character code. \\
Mouse & Pointing device moved on a surface to control the on-screen pointer. \\
Trackball & Pointing device whose ball is rotated in place; the device stays still. \\
Touchpad & Flat, touch-sensitive pad on a laptop; the pointer moves on the screen. \\
Joystick & Lever that moves on two or three axes, used for games and control. \\
Light pen & Pen-shaped input device with a light sensor, used on a CRT screen. \\
Stylus / digitizer & Pen plus tablet that converts strokes into digital coordinates. \\
Scanner & Converts a paper document or picture into a digital image. \\
Biometrics & Input from a physical or behavioural human trait (fingerprint, iris, face). \\
\bottomrule
\end{tabular}
\end{center}

\section{Classifying input devices}
Input devices are grouped by the type of data they capture. This is the
classification JKSSB question-setters use most often, so learn one example
for each row.

\begin{center}
\begin{tabular}{p{0.26\textwidth}p{0.60\textwidth}}
\toprule
\textbf{Category} & \textbf{Typical devices} \\
\midrule
Text input & Keyboard \\
Pointing input & Mouse, trackball, touchpad, joystick, light pen, stylus, touchscreen \\
Audio input & Microphone \\
Image / video input & Scanner, digital camera, webcam \\
Biometric input & Fingerprint, iris, and face sensors \\
Code / tag input & Barcode, QR code, RFID \\
\bottomrule
\end{tabular}
\end{center}

A single device may belong to more than one row if it captures more than one
kind of data, but for exam purposes use the primary role: a touchscreen is a
pointing input device, a webcam is an image/video input device, and a
microphone is an audio input device.

\section{Direct vs indirect input}
A second useful classification is whether the data is captured at its source
or encoded by the user.

\textbf{Direct input devices} capture the data themselves. Examples are the
touchscreen, light pen, scanner, barcode reader, biometric sensor, and
microphone: the device reads the object, code, trait, or sound and turns it
directly into computer data.

\textbf{Indirect input devices} require the user to act on the device to
encode the data. Examples are the keyboard, mouse, trackball, touchpad, and
joystick: the user translates an intention into keystrokes or pointer
movements, which the device then sends to the computer.

\begin{definitionbox}{Direct input}
Input in which the device captures the data at its source (for example, a
scanner reading a page or a biometric sensor reading a fingerprint) without
the user encoding it by hand.
\end{definitionbox}

The distinction matters because it explains why some devices are fast and
error-free while others depend on the user. A barcode reader or a biometric
sensor captures the data exactly as it is, so there is no typing mistake; a
keyboard, being indirect, is only as accurate as the person typing. In a
question, if the stem stresses that the device \emph{reads} or \emph{captures}
an object, code, trait, or sound, think of a direct input device; if it
stresses that the \emph{user types or moves} something to enter the data,
think of an indirect input device.

\section{The keyboard: text input}
The keyboard is the most common text input device. Its standard letter
layout is called \textbf{QWERTY}, after the first six letters of the top row.
A keyboard has alphanumeric keys, function keys (F1 to F12), modifier keys
such as Ctrl, Alt, and Shift, and usually a numeric keypad. Each time a key
is pressed, the keyboard sends the corresponding character code to the
computer. The keyboard is an indirect input device: the user types the data,
and the device encodes it.

Keyboards differ in construction even though their role is the same. A
\textbf{membrane keyboard} uses a flexible membrane under the keys and is
cheap and quiet. A \textbf{mechanical keyboard} uses a separate physical
switch under each key and gives a firmer, more definite feel. A
\textbf{wireless keyboard} connects to the computer by Bluetooth or by a USB
receiver instead of a cable. Whatever the construction, a keyboard is still a
text input device: it sends character codes for the keys that are pressed.

\section{Pointing devices}
A pointing device controls the position of an on-screen pointer (cursor) and
is used to point, click, and drag. The table below fixes each pointing device
and its nearest neighbour, because JKSSB questions often ask candidates to
tell two similar pointing devices apart.

\begin{center}
\begin{tabular}{p{0.16\textwidth}p{0.38\textwidth}p{0.34\textwidth}}
\toprule
\textbf{Device} & \textbf{How it works} & \textbf{Nearest neighbour} \\
\midrule
Mouse & The device slides on a surface and the pointer follows. & Trackball: the ball rotates while the device stays still. \\
Trackball & The ball is rotated in place; the device is stationary. & Mouse: the whole device moves. \\
Touchpad & A flat finger pad moves the pointer on the screen (indirect). & Touchscreen: the display itself is touched (direct). \\
Joystick & A lever moves on two or three axes for games and control. & Mouse: flat-surface pointing. \\
Light pen & A light sensor is pointed at a CRT screen. & Stylus: used on a digitizer surface. \\
Stylus + digitizer & A pen stroke is converted into digital coordinates. & Light pen: uses a light sensor, not a tablet. \\
\bottomrule
\end{tabular}
\end{center}

\subsection{The mouse}
The \textbf{mouse} is a hand-held pointing device moved on a flat surface. It
has one or more buttons and usually a scroll wheel. An \textbf{optical or
laser mouse} tracks movement with light, which is the common modern design;
an older \textbf{mechanical mouse} used a rolling ball. A mouse is an input
device (PYQ-08).

\subsection{The trackball}
The \textbf{trackball} is, in effect, an upside-down mouse: the ball is on
top and is rotated with the fingers while the device itself stays still. This
saves desk space because the device does not have to be moved. The
nearest-neighbour trap is that a trackball and a mouse both point, but with a
mouse the \emph{device} slides, while with a trackball the \emph{ball}
rotates in place.

\subsection{The trackpad or touchpad}
The \textbf{trackpad} (also called the touchpad) is a flat, touch-sensitive
pad found on laptops. The user slides a finger across the pad to move the
pointer on the screen and taps to click. It is an indirect device because the
finger does not touch the object on the screen; the pad only controls the
pointer.

\subsection{The joystick}
The \textbf{joystick} is a lever that can be moved on two or three axes. It is
used mainly for games and for some industrial control tasks.

\subsection{The light pen}
The \textbf{light pen} is a pen-shaped input device with a light sensor at its
tip. It is pointed at a CRT screen to select items or to draw. A light pen is
an input device, not an output device (TRAP-12).

\section{Stylus and digitizer tablet}
A \textbf{stylus} is a pen-like tool used to write or draw on a tablet or a
touch screen. A \textbf{digitizer} (graphics tablet) is a flat surface that
converts the movement of the stylus into digital coordinates. Together they
are used for freehand drawing, capturing signatures, computer-aided design
(CAD), and handwriting recognition. Because a digitizer captures a drawing as
data, it is an input device (TRAP-12). Do not confuse the stylus with the
light pen: the light pen works with a light sensor, while the stylus works
with a digitizer surface.

\section{Touchscreen}
A \textbf{touchscreen} is a display that also works as an input device. The
user touches the screen directly to select and enter data. Because the item
touched is the item selected, a touchscreen is a direct input device. The two
common technologies are \textbf{resistive} and \textbf{capacitive}
touchscreens. Touchscreens are used in smartphones, tablets, ATMs, and
information kiosks.

The nearest-neighbour trap is the touchscreen versus the touchpad. On a
touchscreen the user touches the display itself, so the input is direct. On a
touchpad the user touches a separate pad and the pointer moves on the screen,
so the input is indirect. A touchpad is not a touchscreen.

\section{Scanners, cameras, and webcams}
A \textbf{scanner} converts a paper document or a picture into a digital
image. Common types are the \textbf{flatbed} scanner (the page is placed on a
glass surface), the sheet-fed scanner, the handheld scanner, and the drum
scanner. Software such as OCR (Optical Character Recognition) or OMR (Optical
Mark Recognition) can then read text or marks from the scanned image. A
scanner is an input device (TRAP-12).

A \textbf{digital camera} captures photographs as digital images and stores
them on a memory card. A \textbf{webcam} is a small digital camera connected
to a computer and used mainly for live video, such as video calls and
streaming. Both are image/video input devices.

\section{Barcode, QR code, and RFID}
A \textbf{barcode} is a one-dimensional pattern of parallel lines of varying
width, read by a barcode scanner. It is used on product labels and on book
ISBNs.

A \textbf{QR code} (Quick Response code) is a two-dimensional matrix of black
and white squares. It can store more data than a barcode and can be read from
any direction. The full form of QR is \textbf{Quick Response}.

\textbf{RFID} stands for \textbf{Radio Frequency Identification}. An RFID
system has a tag and a reader; the two exchange radio waves, so no line of
sight is needed, unlike a barcode. RFID is used in toll collection, access
cards, and inventory tracking. A barcode, a QR code, and an RFID tag are all
read by an input device that feeds the data into the computer.

\section{Microphone and biometric input}
A \textbf{microphone} is an audio input device. It converts sound, such as
voice or music, into an electrical signal that the computer digitises. It is
used for recording, voice calls, and speech recognition. The matching output
device is the \textbf{speaker}: a microphone takes audio in, a speaker sends
audio out.

\textbf{Biometric} input identifies a person from a physical or behavioural
trait. Common traits are the \textbf{fingerprint}, the \textbf{iris}, the
\textbf{face}, the voice, and the palm. A sensor captures the trait and feeds
the measured data to the computer. Biometrics are used for attendance
systems, device unlocking, and authentication. Because the sensor reads a
human trait and passes the data to the system, biometric input is a direct
input form.

\begin{example}
A student wants to scan a printed page into the computer. The correct device
is a scanner, which converts the page into a digital image. If the student
instead wanted to type the same text by hand, the device would be a keyboard,
which is an indirect text input device. If the goal were to capture a live
video of the student, the device would be a webcam.
\end{example}

\section{Summary of input devices}
The table below gathers the devices in this lesson for quick revision. Use it
as a checklist: for each device, be ready to state its category and its
one-line function, and to name the device it is most often confused with.

\begin{center}
\begin{tabular}{p{0.24\textwidth}p{0.24\textwidth}p{0.42\textwidth}}
\toprule
\textbf{Device} & \textbf{Category} & \textbf{Function} \\
\midrule
Keyboard & Text & Converts a keystroke into a character code. \\
Mouse & Pointing & Moves the on-screen pointer; point, click, drag. \\
Trackball & Pointing & Ball rotated in place; device stays still. \\
Touchpad & Pointing & Flat finger pad; pointer moves on the screen. \\
Joystick & Pointing & Lever on two or three axes for games/control. \\
Light pen & Pointing & Light sensor pointed at a CRT screen. \\
Stylus + digitizer & Pointing & Converts pen strokes into digital coordinates. \\
Touchscreen & Pointing (direct) & Touch the display to select and enter. \\
Scanner & Image & Converts a document or picture into a digital image. \\
Digital camera & Image/video & Captures photographs as digital images. \\
Webcam & Image/video & Small camera for live video on a computer. \\
Microphone & Audio & Converts sound into an electrical signal. \\
Biometric sensor & Biometric & Reads a fingerprint, iris, or face trait. \\
Barcode / QR code & Code/tag & Read optical codes that identify an item. \\
RFID & Code/tag & Reads a tag by radio waves; no line of sight. \\
\bottomrule
\end{tabular}
\end{center}

\section{Points of confusion between similar devices}
Beyond classifying a device, JKSSB items often test a single distinction
between two devices that look or sound alike. The table below collects the
pairs used in this lesson.

\begin{center}
\begin{tabular}{p{0.30\textwidth}p{0.60\textwidth}}
\toprule
\textbf{Pair} & \textbf{The distinction to state} \\
\midrule
Mouse vs trackball & Mouse: the device slides. Trackball: the ball rotates, device still. \\
Touchscreen vs touchpad & Touchscreen: touch the display (direct). Touchpad: separate pad moves the pointer (indirect). \\
Light pen vs stylus & Light pen: light sensor on a CRT screen. Stylus: used with a digitizer tablet. \\
Scanner vs digital camera & Scanner: converts a flat document or picture into an image. Camera: captures a scene. \\
Barcode vs QR code & Barcode: 1D parallel lines. QR code: 2D matrix, holds more data. \\
Barcode/QR vs RFID & Barcode/QR are read optically and need line of sight; RFID uses radio waves. \\
Microphone vs speaker & Microphone: audio input. Speaker: audio output. \\
Data glove vs haptic glove & Data glove: captures hand movement (input). Haptic glove: gives force feedback (output). \\
\bottomrule
\end{tabular}
\end{center}

\section{Exam retrieval and common confusions}
Six distinctions prevent most errors on this topic.
\begin{enumerate}
  \item A light pen is an \textbf{input} device; a plotter is an \textbf{output}
    device. Do not swap them (TRAP-12).
  \item A punch-card reader, a digitizer tablet, and a scanner are all
    \textbf{input} devices. A drum printer and a sound card are
    \textbf{output} devices (TRAP-12).
  \item A \textbf{trackball} rotates its ball in place; a \textbf{mouse}
    slides over a surface.
  \item A \textbf{touchscreen} is touched directly; a \textbf{touchpad} is a
    separate pad that moves the pointer indirectly.
  \item \textbf{RFID} uses radio waves and needs no line of sight; a
    \textbf{barcode} and a \textbf{QR code} are read optically.
  \item \textbf{Haptic gloves}, a \textbf{tactile display}, and a
    \textbf{Voice Response Unit (VRU)} give feedback to the user and are
    \textbf{output}; a \textbf{data glove} captures hand movement and is
    \textbf{input} (PYQ-09).
\end{enumerate}

\begin{example}
An item lists four devices --- light pen, punch-card reader, plotter, and
scanner --- and asks which are input devices. Classify each one: the light
pen is input, the punch-card reader is input, the plotter is output, and the
scanner is input. The answer is the light pen, the punch-card reader, and the
scanner. The plotter is the only output device in the list.
\end{example}

\subsection*{Exercises}
\begin{enumerate}[label=\arabic*.]
  \item Define an input device and give three examples.
  \item Classify these devices by the type of data they capture: keyboard,
    microphone, scanner, fingerprint sensor.
  \item Distinguish direct input from indirect input, giving one example of
    each.
  \item Explain the difference between a trackball and a mouse, and between a
    touchscreen and a touchpad.
  \item Differentiate a barcode, a QR code, and RFID.
  \item State whether each is input or output: light pen, plotter, digitizer
    tablet, drum printer, sound card.
\end{enumerate}

\subsection*{Solutions}
\begingroup\small
\begin{enumerate}[label=\arabic*.]
  \item An input device accepts data or control signals from the user or the
    environment and passes them to the CPU in machine-readable form.
    Examples: keyboard, mouse, scanner (any three valid devices).
  \item Keyboard --- text input; microphone --- audio input; scanner ---
    image input; fingerprint sensor --- biometric input.
  \item Direct input captures the data at its source, such as a scanner or a
    biometric sensor. Indirect input needs the user to encode the data, such
    as a keyboard or a mouse.
  \item With a mouse, the device slides over a surface; with a trackball, the
    ball is rotated in place and the device stays still. A touchscreen is
    touched directly on the display; a touchpad is a separate pad that moves
    the pointer indirectly.
  \item A barcode is a 1D pattern of lines read optically; a QR code is a 2D
    matrix that holds more data; RFID uses radio waves between a tag and a
    reader and needs no line of sight.
  \item Light pen --- input; plotter --- output; digitizer tablet --- input;
    drum printer --- output; sound card --- output.
\end{enumerate}
\endgroup
\end{document}
